\section{MIN's schedules and the machine, in detail}\label{app:more}
\paragraph{The fewest-admission schedule, host by host.} \Cref{tab:job104,tab:job105} list jobs 104 and 105: MIN's set
loaded by the CPU or in the step, as the greedy rule picks it or as the schedule with the fewest admissions does. On the
engine's 20 problems the fewest-admission schedule copies \jfCopyRatioMin--\jfCopyRatioMax{} times as many experts in
the step as the greedy one, for \jfMissDevMin--\jfMissDevMax\% more misses (more than its replay predicts; its measured
gains include that cost). On every launch it gains over the deployed cache with either load and at both budgets, each
95\% interval above 1; loaded by the CPU on the machine with the fastest link, it is \jhFastBypassMidBehind{} behind the
greedy set at 25\%.

% generated by scripts/job103.py
\begin{table}[t]\centering\footnotesize
\caption{Which of MIN's hit-optimal sets the oracle follows (job 104): speed relative to the deployed cache on the same machine, with MIN's set loaded by the CPU (two reads: \MinTwo, \FewTwo) or in the step (one read: \MinOne, \FewOne), as the greedy rule picks it or as the schedule with the fewest admissions does; and the copies per token each makes in the step. Machines sorted by the probe's link-to-CPU ratio.}\label{tab:job104}
\setlength\tabcolsep{2.5pt}\resizebox{\linewidth}{!}{%
\begin{tabular}{@{}lrlrrrrrr@{}}\toprule
 & Link/ & & \multicolumn{2}{c}{By the CPU} & \multicolumn{2}{c}{In the step} & \multicolumn{2}{c}{Copies/token} \\
Machine & CPU & Budget & \MinTwo & \FewTwo & \MinOne & \FewOne & \MinOne & \FewOne \\\midrule
Pf again & 0.29 & gpt-oss 11\% & 0.94 & 1.11 & 0.86 & 1.04 & 19.8 & 11.8 \\
Pf again & 0.29 & gpt-oss 25\% & 1.13 & 1.22 & 0.96 & 1.06 & 9.7 & 7.0 \\
285K, second & 0.49 & gpt-oss 11\% & 1.03 & 1.08 & 1.05 & 1.18 & 19.8 & 11.8 \\
285K, second & 0.49 & gpt-oss 25\% & 1.20 & 1.21 & 1.12 & 1.17 & 9.7 & 7.0 \\
Pg again (5950X) & 0.63 & gpt-oss 11\% & 1.02 & 1.06 & 1.26 & 1.34 & 19.8 & 11.8 \\
Pg again (5950X) & 0.63 & gpt-oss 25\% & 1.20 & 1.15 & 1.32 & 1.35 & 9.7 & 7.0 \\
\bottomrule\end{tabular}}\end{table}

% generated by scripts/job105.py
\begin{table}[t]\centering\footnotesize
\caption{Job 105, registered before the runs, nothing fitted: the deployed cache's time per token against the plain form of \cref{eq:sum} (without the overlap term, which came after this job), the GPU's compute profiled on the same host before the timed runs plus the run's counted host reads at the machine's best rate; and the speed, relative to the deployed cache, of the deployed cache with the engine's layer-ahead copy (each layer's top predicted expert of the next layer copied ahead of use) and of MIN's greedy and fewest-admission sets copied in the step. Hosts sorted by the probe's link-to-CPU ratio.}\label{tab:job105}
\setlength\tabcolsep{2.5pt}\resizebox{\linewidth}{!}{%
\begin{tabular}{@{}lrlrrrrrr@{}}\toprule
 & Link/ & & $G_{\text{prof}}$ & \multicolumn{2}{c}{Deployed (ms)} & Layer- & \multicolumn{2}{c}{MIN, in step} \\
Host & CPU & Budget & (ms) & plain Eq.~3 & measured & ahead & greedy & fewest \\\midrule
Pf again & 0.28 & gpt-oss 11\% & 4.1 & 13.6 & 13.9 & 0.72 & 0.86 & 1.03 \\
Pf again & 0.28 & gpt-oss 25\% & 4.4 & 9.8 & 9.5 & 0.69 & 0.97 & 1.06 \\
Threadripper 9960X & 0.32 & gpt-oss 11\% & 4.1 & 9.3 & 9.1 & 0.91 & 0.91 & 1.04 \\
Threadripper 9960X & 0.32 & gpt-oss 25\% & 4.4 & 7.3 & 6.8 & 0.87 & 0.98 & 1.04 \\
EPYC 7402P & 0.45 & gpt-oss 11\% & 4.3 & 19.1 & 18.2 & 0.81 & 1.01 & 1.15 \\
EPYC 7402P & 0.45 & gpt-oss 25\% & 4.5 & 12.9 & 11.8 & 0.78 & 1.09 & 1.15 \\
O4 again (9950X) & 1.04 & gpt-oss 11\% & 4.4 & 20.8 & 20.5 & 1.04 & 1.36 & 1.34 \\
O4 again (9950X) & 1.04 & gpt-oss 25\% & 4.4 & 13.2 & 12.4 & 1.01 & 1.42 & 1.39 \\
\bottomrule\end{tabular}}\end{table}

% generated by scripts/job106.py
\begin{table}[t]\centering\footnotesize
\caption{Job 106, registered before the runs, every host ($^*$: deployed cache varied by more than the registered 2\% between rounds): the deployed cache's time per token against \cref{eq:sum} (GPU compute profiled on the same host before the timed runs; the run's own counters), and the speed relative to the deployed cache of the two online admission rules (\emph{dk}: the decayed count with a larger admission margin; \emph{lrn}: the learned reuse predictor with a margin; both chosen on other text) and of MIN's fewest-admission set copied in the step, with round-level 95\% intervals (rounds, then problems: three rounds at gpt-oss 11\%, two at 25\%; Qwen3 ran one round, so its intervals are over problems). Hosts sorted by the probe's link-to-CPU ratio.}\label{tab:job106all}
\setlength\tabcolsep{2.2pt}\resizebox{\linewidth}{!}{%
\begin{tabular}{@{}lrlrrrlll@{}}\toprule
 & Link/ & & $G$ & \multicolumn{2}{c}{Deployed (ms)} & & & MIN, fewest \\
Host & CPU & Budget & (ms) & Eq.~3 & meas. & dk & lrn & in the step \\\midrule
EPYC 7302$^*$ & 0.14 & gpt-oss 11\% & 4.3 & 12.2 & 17.7 & 1.06 [1.01, 1.10] & 0.97 [0.92, 1.01] & 0.85 [0.82, 0.88] \\
EPYC 7302$^*$ & 0.14 & gpt-oss 25\% & 4.7 & 8.7 & 12.2 & 1.08 [1.00, 1.17] & 1.10 [1.06, 1.14] & -- \\
EPYC 7302$^*$ & 0.14 & Qwen3 12.5\% & 5.8 & 20.4 & 25.4 & 1.07 [1.02, 1.13] & 0.93 [0.86, 1.00] & -- \\
EPYC 7302$^*$ & 0.14 & Qwen3 25\% & 6.3 & 14.3 & 17.2 & 1.01 [0.93, 1.08] & 0.84 [0.76, 0.94] & -- \\
Xeon 8347C$^*$ & 0.23 & gpt-oss 11\% & 4.1 & 12.0 & 19.5 & 1.28 [0.99, 1.56] & 1.08 [0.80, 1.45] & 1.33 [1.00, 1.61] \\
Xeon 8347C$^*$ & 0.23 & gpt-oss 25\% & 4.5 & 8.6 & 10.1 & 1.06 [1.05, 1.08] & 1.02 [0.99, 1.04] & -- \\
Xeon 8347C$^*$ & 0.23 & Qwen3 12.5\% & 5.3 & 21.4 & 45.0 & 1.96 [1.67, 2.24] & 1.29 [0.94, 1.74] & -- \\
Xeon 8347C$^*$ & 0.23 & Qwen3 25\% & 5.8 & 14.9 & 16.6 & 0.86 [0.71, 1.02] & 0.73 [0.59, 0.88] & -- \\
Pf again (285K) & 0.29 & gpt-oss 11\% & 4.3 & 13.8 & 13.7 & 1.02 [1.01, 1.04] & 1.03 [1.02, 1.03] & 1.02 [1.01, 1.03] \\
Pf again (285K) & 0.29 & gpt-oss 25\% & 4.7 & 9.6 & 9.5 & 1.06 [1.04, 1.07] & 1.05 [1.04, 1.07] & -- \\
Pf again (285K) & 0.29 & Qwen3 12.5\% & 5.4 & 24.8 & 24.4 & 1.02 [1.00, 1.03] & 1.02 [1.01, 1.03] & -- \\
Pf again (285K) & 0.29 & Qwen3 25\% & 5.9 & 16.9 & 16.2 & 1.04 [1.03, 1.04] & 1.03 [1.02, 1.03] & -- \\
Threadripper 9960X again & 0.32 & gpt-oss 11\% & 4.1 & 9.1 & 9.4 & 1.00 [0.99, 1.02] & 0.98 [0.97, 1.00] & 1.06 [1.05, 1.06] \\
Threadripper 9960X again & 0.32 & gpt-oss 25\% & 4.4 & 6.9 & 7.0 & 1.02 [1.00, 1.04] & 0.99 [0.97, 1.00] & -- \\
Threadripper 9960X again & 0.32 & Qwen3 12.5\% & 5.4 & 15.1 & 14.4 & 1.02 [1.01, 1.02] & 0.99 [0.99, 1.00] & -- \\
Threadripper 9960X again & 0.32 & Qwen3 25\% & 5.7 & 11.2 & 10.5 & 1.01 [1.01, 1.01] & 0.98 [0.98, 0.99] & -- \\
Pe again (9950X) & 1.02 & gpt-oss 11\% & 4.5 & 20.4 & 20.5 & 1.01 [1.01, 1.01] & 1.01 [1.01, 1.02] & 1.34 [1.33, 1.35] \\
Pe again (9950X) & 1.02 & gpt-oss 25\% & 4.5 & 12.8 & 12.4 & 1.01 [1.01, 1.01] & 1.01 [1.00, 1.02] & -- \\
Pe again (9950X) & 1.02 & Qwen3 12.5\% & 5.8 & 35.0 & 36.7 & 1.01 [1.01, 1.01] & 1.02 [1.02, 1.02] & -- \\
Pe again (9950X) & 1.02 & Qwen3 25\% & 5.9 & 22.2 & 22.5 & 1.01 [1.01, 1.01] & 1.01 [1.01, 1.02] & -- \\
\bottomrule\end{tabular}}\end{table}


\paragraph{Job 106, every host.} \Cref{tab:job106all} adds the two server machines new to job 106. On the
\jiNllBadName{} host the engine computed wrong outputs: its teacher-forced loss on the deployed cache was
\jiNllBadMin--\jiNllBadMax{} nats per token across rounds and models, against \jiNllGood{} (gpt-oss) and \jiNllGoodQ{}
(Qwen3) on every other host and in its own untimed routing pass, and its counters changed between rounds. On the Xeon
8347C host the deployed cache varied by \jiUnstableVarMax\% between rounds at gpt-oss 11\%, against at most 2\%
registered, and the greedy copy gained in two of its three rounds where we had predicted it would lose. When their jobs
reached the gate, before any download, \texttt{free} reported \jiUsedUnstMin--\jiUsedUnstMax\,GB of host memory in use
(\jiUsedStableMin--\jiUsedStableMax\,GB on the other hosts), and their probes' CPU read rates fell when more threads
read than the container had usable physical cores. On both, \cref{eq:sum} under-predicted the deployed cache by
\jiUnstErrMin--\jiUnstErrMax\%. We report their clauses (in the supplement) and draw nothing from them about the relation on those runs; both are servers, and job 107's one valid server was under-predicted as well (\cref{sec:gap}). The three
machines that ran stably were all rented before: Pf, the Threadripper 9960X of job 105, and panel host Pe (same GPU).

% generated by scripts/job107.py
\begin{table}[t]\centering\footnotesize
\caption{Job 107, registered before the runs, on the machines never rented before that passed the registered gate before any download ($^\dagger$: rounds further apart than the registered 2\%; such a machine stopped after its gpt-oss 11\% rounds and is outside the predictions): the deployed cache's time per token against \cref{eq:sum} (GPU compute profiled on the same machine before the timed runs; the run's own counters), and the speed relative to the deployed cache of \Margin{} and of MIN's fewest-admission set copied in the step or loaded by the CPU, with round-level 95\% intervals (rounds, then problems: two rounds of gpt-oss, one of Qwen3, whose intervals are over problems). Machines sorted by the probe's link-to-CPU ratio.}\label{tab:job107}
\setlength\tabcolsep{2.2pt}\resizebox{\linewidth}{!}{%
\begin{tabular}{@{}lrlrrrlll@{}}\toprule
 & Link/ & & $G$ & \multicolumn{2}{c}{Deployed (ms)} & & \multicolumn{2}{c}{\FewOne, \FewTwo} \\
Machine & CPU & Budget & (ms) & Eq.~\ref{eq:sum} & meas. & \Margin & in the step & by the CPU \\\midrule
EPYC 7663$^{\dagger}$ & 0.21 & gpt-oss 11\% & 4.3 & 11.3 & 14.3 & 1.03 [1.00, 1.05] & 0.97 [0.94, 1.01] & 1.13 [1.07, 1.18] \\
AMD engineering sample & 0.44 & gpt-oss 11\% & 4.1 & 8.2 & 20.0 & 1.05 [1.04, 1.06] & 1.16 [1.14, 1.17] & 1.13 [1.11, 1.14] \\
AMD engineering sample & 0.44 & gpt-oss 25\% & 4.2 & 6.4 & 13.4 & 0.99 [0.98, 1.00] & -- & -- \\
AMD engineering sample & 0.44 & Qwen3 12.5\% & 5.6 & 13.6 & 32.7 & 1.03 [1.02, 1.04] & -- & -- \\
AMD engineering sample & 0.44 & Qwen3 25\% & 5.9 & 10.4 & 23.7 & 1.07 [1.06, 1.09] & -- & -- \\
EPYC 9754$^{\dagger}$ & 0.51 & gpt-oss 11\% & 4.2 & 10.1 & 14.7 & 0.97 [0.95, 0.99] & 1.18 [1.15, 1.20] & 1.05 [1.02, 1.08] \\
Ryzen 9 5900XT & 0.75 & gpt-oss 11\% & 4.4 & 27.4 & 29.5 & 1.02 [1.02, 1.02] & 1.42 [1.41, 1.43] & 1.04 [1.04, 1.04] \\
Ryzen 9 5900XT & 0.75 & gpt-oss 25\% & 4.5 & 16.6 & 17.0 & 1.02 [1.01, 1.02] & -- & -- \\
Ryzen 9 5900XT & 0.75 & Qwen3 12.5\% & 5.8 & 48.3 & 54.4 & 1.02 [1.02, 1.02] & -- & -- \\
Ryzen 9 5900XT & 0.75 & Qwen3 25\% & 5.9 & 29.7 & 31.7 & 1.02 [1.02, 1.02] & -- & -- \\
\bottomrule\end{tabular}}\end{table}


\paragraph{Job 107, every host.} \Cref{tab:job107} lists the four machines that passed the gate before any download.
The EPYC 9754 and the EPYC 7663 ran their gpt-oss 11\% rounds \jlInvalidSpreadMin\% and \jlInvalidSpreadMax\% apart,
beyond the registered 2\%, and stopped there, as registered; the relation under-predicted both, by \jlInvErrMin--\jlInvErrMax\%. Three
more machines stopped at the gate with \jlUsedGateMin--\jlUsedGateMax\,GB of host memory in use. The registered
replacement list ran out after the first gate failure and two hosts that failed the round check; the two hosts added
after that were chosen by a rule written into the job script before each started (cheapest verified offer not rented
before, with fast downloads and cheap traffic), and both stopped at the gate. On the two valid machines, the relation
predicted the admission margin's speed from its own counters within a median \jlDkPredMed{} (up to \jlDkPredMax, on the
server), against \jlDkNoChangeMed{} for predicting no change. The profiled GPU compute was \jlGMin--\jlGMax\,ms on
gpt-oss and \jlGQMin--\jlGQMax\,ms on Qwen3, as on every earlier host: the GPU's own work does not depend on the CPU.

\paragraph{The machine decides how much.} Across machines, the gain of every state that moves reads from the CPU to
the link follows the link-to-CPU ratio (Spearman over the machines that ran it through job 104, 95\% bootstrap
interval): MIN copied in the step \rcFetchRatio{} [\rcFetchRatioLo, \rcFetchRatioHi] (\rcFetchRatioN{} machines), copied
ahead \rcPacedRatio{} [\rcPacedRatioLo, \rcPacedRatioHi] (\rcPacedRatioN), admitting every miss \rcAaRatio{}
[\rcAaRatioLo, \rcAaRatioHi] (\rcAaRatioN). MIN loaded by the CPU, which reads every admission twice and copies it in
the background, tracks the link's absolute rate more closely (\rcBypassBp{} [\rcBypassBpLo, \rcBypassBpHi],
\rcBypassBpN{} machines, against \rcBypassRatio{} for the ratio; the two intervals overlap, and the feature was picked
from a search over all probe features). The greedy schedule copied in the step gains
\jeNewFetchLowMin--\jeNewFetchLowMax$\times$ at 11\% on \jeNewN{} further machines with ratios \jeNewRatioMin{} and
\jeNewRatioMax{}, and loses only on the machines whose link reads at under a third of their CPU rate; made ahead of use,
the same copies lose beyond their interval on no machine they ran on (at worst \jbPfPaced$\times$, on Pf).

\paragraph{Machines differ more than problems.} Problems rank largely the same on every machine: Kendall's coefficient
of concordance of MIN-copied-in-the-step's per-problem gain over \ppMachines{} machines and \ppProblems{} problems is
\ppWFetchLow{} at 11\% and \ppWFetchMid{} at 25\%, and a problem's gain follows how far the deployed policy's reads
exceed MIN's on its trace (Spearman \ppRhoFetchLow{} at 11\%, \ppRhoFetchMid{} at 25\%). Yet machines carry
\ppVarHostsFetchMax\% of the variance of its log gain and problems \ppVarProbFetchMin--\ppVarProbFetchMax\%.
